\documentclass[uplatex,dvipdfmx]{jsarticle}
\usepackage{amsmath,amssymb,amsthm}
\usepackage{enumitem}
\usepackage{seqsplit}
\usepackage[
  dvipdfmx,
  unicode,
  bookmarks=true,
  bookmarksnumbered=true,
  bookmarksopen=true,
  bookmarksdepth=2
]{hyperref}
% pxjahyper converts Japanese strings in PDF metadata and outlines to Unicode.
% It must be loaded after hyperref when compiling with upLaTeX + dvipdfmx.
% TeX Live 2022's pxjahyper uses the deprecated spelling of one optional OTF
% package hook.  Silence only that kernel message while loading the package.
\ExplSyntaxOn
\msg_redirect_name:nnn { hooks } { generic-deprecated } { none }
\ExplSyntaxOff
\usepackage[nootfcid,nootfmacros]{pxjahyper}
\ExplSyntaxOn
\msg_redirect_name:nnn { hooks } { generic-deprecated } { warning }
\ExplSyntaxOff
\usepackage{bookmark}
\hypersetup{
  hidelinks,
  pdftitle={有理関数体上の Witt 型 Local--Global 原理},
  pdfauthor={selpo},
  pdfsubject={Lean 形式化と対応する完全な数学的証明}
}

\newcommand{\RF}{R(X)}
\newcommand{\RatFunc}{\operatorname{RatFunc}}
\newcommand{\lean}[1]{\texttt{\expandafter\seqsplit\expandafter{\detokenize{#1}}}}

\newtheorem{theorem}{定理}
\newtheorem{proposition}[theorem]{命題}
\newtheorem{lemma}[theorem]{補題}
\newtheorem{corollary}[theorem]{系}
\theoremstyle{definition}
\newtheorem{definition}[theorem]{定義}
\theoremstyle{remark}
\newtheorem{remark}[theorem]{注意}
\renewcommand{\proofname}{証明}

\title{有理関数体上の Witt 型 Local--Global 原理}
\author{selpo}
\date{}

\begin{document}
\pdfbookmark[0]{有理関数体上の Witt 型 Local--Global 原理}{document-title}
\maketitle

\begin{abstract}
$R$ を実閉体、$F=R(X)$ を一変数有理関数体とする。
Lean プロジェクト \texttt{lean-ratfunc-witt-}\allowbreak\texttt{local-global} で
形式化した次の定理の証明を与える：
次元3以上の有限次元二次形式は、$F$ 上 isotropic であることと、$F$ の任意の
実閉順序拡大上 isotropic であることが同値である。三項基底定理のための ordered real
closure を構成し、その後、符号分離を行う binary--tail compression によって次元帰納する。
最後に数学上の補題と Lean 宣言の対応を記録する。
\end{abstract}

\section{主定理と記法}

以下 $R$ は実閉体、$F=R(X)$ とする。有限集合 $I$ と係数族
$a=(a_i)_{i\in I}$ に対し
\[
  \langle a\rangle(x)=\sum_{i\in I}a_i x_i^2
\]
と書く。非零ベクトル $x\in F^I$ でこの和が零になるとき、この対角形式を
\emph{isotropic} と呼ぶ。regular な対角形式について、$F$ の ordering において
全係数がstrictに正、または全係数がstrictに負になるとき、その ordering で形式は
definite であるという。

\begin{theorem}[対角形式の local--global 定理]\label{thm:diagonal-ja}
$I$ を有限集合、$|I|\ge3$、$a:I\to F$ とする。このとき次は同値である。
\begin{enumerate}[label=(\roman*)]
  \item $\langle a\rangle$ は $F$ 上 isotropic である。
  \item 全係数が同じ strict sign を持つような $F$ の ordering は存在しない。
  \item 任意の実閉 ordered $F$-algebra $K$ 上で、scalar extension
        $\langle a\rangle_K$ は isotropic である。
\end{enumerate}
Lean の same-universe 版では (iii) の $K$ を固定した universe 内で量化する。
別 universe の版では、量化対象が空になることを避けるため、対応する ordered real
closure の存在を明示的な引数に取る。
\end{theorem}

\begin{theorem}[一般二次形式の local--global 定理]\label{thm:qf-ja}
$V$ を $\dim_F V\ge3$ の有限次元 $F$-vector space、$q:V\to F$ を
二次形式とする。このとき $q$ が anisotropic でないことと、任意の実閉 ordered
extension 上の scalar extension が anisotropic でないことは同値である。
形式化された定理は対角形式の場合と同様に same-universe の対象を量化し、別 universe
の版では対応する ordered real closure を明示的に仮定する。
\end{theorem}

係数が零なら対応する座標ベクトルが直ちに isotropic である。したがって証明途中では
全係数非零と仮定し、最後に零係数の場合を戻せばよい。

\paragraph{読み方と前提知識}
体、$R[X]$ の一意分解、Chinese remainder theorem、有限次元線形代数、および
Zorn の補題を用いる。それ以外の本証明に
固有な道具は、初めて必要になる箇所で定義し証明する。第2節では ordering を実閉な
検査体へ移すための ordered real closure を構成する。第3節では三項対角形式を
quaternion conic に翻訳し、Lean 本線と同じ polynomial descent を実行する。
第4--5節では、全 ordering で所要の符号を持つ一つの有理関数を構成し、次元を一つ
下げる。第6節で対角形式から一般二次形式へ移る。最後の対応表から、各数学的段階を
実装する Lean 宣言を確認できるので、本文とソースを双方向に読み進められる。

\section{Ordering と実閉拡大}

\begin{lemma}[Definiteness obstruction]\label{lem:definite-ja}
順序体上の対角形式が isotropic なら、その全係数が strict に正となることも、
全係数が strict に負となることもない。
\end{lemma}
\begin{proof}
全係数が正なら、非零ベクトルにおける各項 $a_ix_i^2$ は非負で、少なくとも一項は
正であるから総和は正である。全係数が負の場合は形式全体を負にして同様である。
\end{proof}

\begin{lemma}[実閉体上の対角形式]\label{lem:rcdiag-ja}
$K$ を実閉体、$b_i\in K^\times$ とする。$\langle b\rangle$ が
anisotropic であることと、全ての $b_i$ が同じ strict sign を持つことは同値である。
\end{lemma}
\begin{proof}
一方向は補題~\ref{lem:definite-ja}。異符号の係数 $b_i>0>b_j$ があれば
$-b_j/b_i>0$ であり、実閉体では正元は平方である。$u^2=-b_j/b_i$ として
$i$ 座標を $u$、$j$ 座標を $1$、他を零にすれば isotropic vector を得る。
\end{proof}

形式化では ordered real closure の存在を外部仮定にせず、同じ universe 内で構成する。
ここでは、奇数次拡大に用いる polynomial induction も含め、Lean の証明と同じ構造を
詳しく示す。

\subsection{Ordering を延長する二つの補題}

体 $L$ の \emph{preordering} とは、$0$ と全ての平方を含み、加法と乗法で閉じた
部分集合である。$-1$ を含まないとき \emph{proper} という。
$P\cup(-P)=L$ を満たす proper preordering は ordering と同じものである。

\begin{lemma}[Artin--Schreier extension lemma]\label{lem:preorder-extension-ja}
体の任意の proper preordering は ordering に含まれる。
\end{lemma}
\begin{proof}
与えられた preordering を含む proper preorderings を包含で順序付ける。chain の合併も
proper なので、Zorn の補題により最大元 $P$ がある。$P$ が total であることを示す。
$x,-x\notin P$ と仮定する。最大性より $P\cup\{x\}$ と $P\cup\{-x\}$ が生成する
preordering はいずれも $-1$ を含む。$x$ の偶数冪はすでに $P$ にあるので
\[
  -1=p+xq,\qquad -1=r-xs
  \quad(p,q,r,s\in P)
\]
と書ける。従って $xq=-(1+p)$、$xs=1+r$ であり
\[
  x^2qs=-(1+p)(1+r)\in P.
\]
ここへ $p+r+pr\in P$ を加えると $-1\in P$ となり矛盾する。従って各 $x$ について
$x\in P$ または $-x\in P$ であり、$P$ は ordering である。
\end{proof}

\begin{lemma}[正な射影による判定]\label{lem:projection-ja}
$L/E$ を順序体 $E$ の体拡大とする。$E$-linear map $\lambda:L\to E$ が
\[
  z\ne0\quad\Longrightarrow\quad\lambda(z^2)>0
\]
を満たすなら、$E$ の ordering は $L$ へ延長できる。
\end{lemma}
\begin{proof}
$E$ の非負元と $L$ の平方から生成される preordering $P$ を考える。
$P$ の元は有限和 $\sum_i a_i z_i^2$（$a_i\ge0$）である。もしこの和が $-1$ なら、
$\lambda$ を作用させた右辺は非負である。一方、仮定を $1\ne0$ に適用すると
$\lambda(1)>0$ だから
\[
  \lambda(-1)=-\lambda(1)<0
\]
となり矛盾する。従って $P$ は proper であり、
補題~\ref{lem:preorder-extension-ja}を適用できる。得られる cone は $E$ の非負元を
含む。$a>0$ なのに $-a$ も新しい cone に入れば、$E$ の正元 $a^{-1}$ を掛けて
$-1$ が cone に入るので矛盾する。従って $E$ の正元は正のまま保たれ、これは
もとの ordering の延長である。
\end{proof}

\subsection{奇数次 power basis}

「奇数次拡大は orderable」という部分の偶奇を明確にするため、Lean で用いた
polynomial lemma を証明する。

\begin{lemma}[奇数次 power-basis lemma]\label{lem:odd-power-basis-ja}
$E$ を順序体、$L=E(\theta)$ が奇数次元 $n$ の power basis を持つとする。
$E$ の非負元と $L$ の平方から生成される preordering は proper である。
\end{lemma}
\begin{proof}
$f$ を $\theta$ の monic irreducible polynomial とすると $\deg f=n$ である。
$m\ge1$ に対し
\[
  g=\sum_\nu a_\nu h_\nu^2,\qquad
  a_\nu\ge0,\quad \deg h_\nu<m
\]
と書ける polynomial 全体を $\mathcal C_m$ とする。非零 $g\in\mathcal C_m$ では
最大次数の項の leading coefficient は非負な平方の和であり、相殺しない。従って
\[
  \operatorname{lc}(g)>0,\qquad
  \deg g\text{ は偶数},\qquad
  \deg g<2m.                                           \tag{1}
\]

奇数 $m$ に関する強帰納法で、次数 $m$ の任意の monic irreducible $r$ と
$g\in\mathcal C_m$ に対し
\[
  r\nmid g+1                                           \tag{2}
\]
を示す。反対に $g+1=rk$ とする。$\mathcal C_m$ の元は $-1$ でないので
$g+1\ne0$ である。正の次数を持つ $r$ が $g+1$ を割るため、$g+1$ は非定数である。
従って $g$ も非定数であり、$\deg(g+1)=\deg g$ となる。(1) から
\[
  \deg g=m+\deg k<2m.
\]
$\deg g$ は偶数、$m$ は奇数なので $\deg k$ は奇数で、しかも $\deg k<m$ である。
$k$ の既約因子のうち一つは奇数次数である。それを monic にして $r'$ とする。
$E'=E[T]/(r')$ で評価すると $k=0$、従って $g=-1$ となる。
一方、$\mathcal C_m$ の評価は、$E$ の非負元を係数とする平方の additive span に入る。

$E'$ の各元を次数 $<\deg r'$ の polynomial で表し直すと、この $-1$ の表示を
$g'\in\mathcal C_{\deg r'}$ へ持ち上げられる。すると $r'\mid g'+1$ となり、
$\deg r'<m$ に対する帰納法の仮定に反する。これで (2) が示された。

$L$ の preordering が $-1$ を含むなら、有限個の平方を power basis で表すことにより
$g\in\mathcal C_n$ かつ $g(\theta)=-1$ となる。すると最小多項式 $f$ が $g+1$ を割り、
(2) に反する。
\end{proof}

\begin{proposition}[Ordered real closure]\label{prop:orc-ja}
任意の順序体 $E$ は、その順序を延長する代数的な実閉順序拡大を持つ。
\end{proposition}
\begin{proof}
代数閉包 $\Omega/E$ を固定する。$E$ を含む中間体と、$E$ の順序を延長する
ordering の組を、compatible な包含で順序付ける。chain の合併にも compatible ordering
が入るので、Zorn の補題により最大元 $M$ を得る。

まず $0<a\in M$ とする。$\alpha^2=a$ なる $\alpha\in\Omega$ を取る。
$a$ が $M$ で平方でないなら $T^2-a$ は既約であり、$L=M(\alpha)$ は真の二次拡大
である。$z=x+y\alpha$ の実部への射影は
\[
  \pi(z^2)=x^2+a y^2>0\qquad(z\ne0)
\]
を満たす。補題~\ref{lem:projection-ja}により $L$ を $M$ 上 orderable にできるので、
最大性に反する。よって $M$ の全非負元は平方である。

次に $p\in M[T]$ の次数が奇数とする。$p$ は奇数次数の monic irreducible factor
$r$ を持つ。実際、既約因子の次数の和が奇数なので、そのうち一つは奇数次数である。
$\alpha\in\Omega$ をその根とすれば、$M(\alpha)/M$ は奇数次の power basis を持つ。
補題~\ref{lem:odd-power-basis-ja}と補題~\ref{lem:preorder-extension-ja}により
$M(\alpha)$ は $M$ 上 orderable であり、最大性から $\alpha\in M$ となる。
ゆえに全奇数次多項式は $M$ に根を持つ。

「全非負元が平方」かつ「全奇数次多項式が根を持つ」という順序体の特徴付けにより
$M$ は実閉体である。
\end{proof}

\section{三項基底定理}

\begin{theorem}[三項 local--global 定理]\label{thm:ternary-ja}
$a,b,c\in F^\times$ とする。$\langle a,b,c\rangle$ が $F$ 上 isotropic
であることと、任意の実閉 ordered extension 上 isotropic であることは同値である。
\end{theorem}

以下では Lean の本線で実際に用いる代数的核心を示す。これは quaternion symbol の
descent として整理した polynomial Legendre argument である。

\subsection{Quaternion symbol と polynomial normalization}

$u,v\in F^\times$ に対する quaternion symbol を $(u,v)$ と書く。必要なのは次の
conic による特徴付けだけである：
\[
 (u,v)\text{ が split}
 \quad\Longleftrightarrow\quad
 \exists(x,y,z)\ne(0,0,0),\quad x^2-u y^2-v z^2=0.       \tag{3}
\]
従って
\[
 \langle a,b,c\rangle\text{ が isotropic}
 \quad\Longleftrightarrow\quad
 \left(-\frac ba,-\frac ca\right)\text{ が split}.      \tag{4}
\]
各 slot を非零平方倍しても split 性は変わらない。また
\[
 (u,v)\simeq (u,v(x^2-u y^2))                            \tag{5}
\]
という norm move を、$x^2-u y^2\ne0$ のときに使う。その計算は
\[
 (x^2-u y^2)(X^2-uY^2)
   =(xX+u yY)^2-u(yX+xY)^2
\]
である。(3) の第三座標が非零ならその平方で割ってこの恒等式を適用する。第三座標が
零なら $u$ 自身が平方なので両symbolはsplitする。逆normに同じ計算を使えば逆向きも
従う。以下で用いる quaternion identity は、(5)、slot の
対称性、ternary form の cyclic permutation だけである。

分母を払い、形式全体を一つの共通非零scalar倍し、各係数から平方を除くと、
非零squarefree polynomials $A,B,C$ を得る。これらは pairwise coprime に選べる。
実際、各既約 polynomial $\pi$ について三係数の valuation parity だけを考える。
各平方を除くと parity は $0,1$ であり、形式全体を $\pi$ 倍すると三つの bit が全て
反転する。そこで $1$ の個数が高々一つになる方を選べばよい。これらの操作は
isotropy と「三係数が同じ strict sign」という性質を保つ。

\[
  \mathcal Q(A,B,C)=(-AB,-AC)
\]
と置く。これは (4) の normalized symbol と平方類を除いて同じである。

\subsection{有限 place の合同条件}

有限点と無限遠の符号は Laurent series で実現する。ordered field
$R((\varepsilon))$ では、非零seriesの最初の非零係数が正であるとき、そのseriesを正と
定める。
\[
  X\longmapsto r+\varepsilon,\qquad X\longmapsto r-\varepsilon
\]
は $r$ の直後・直前のorderingを与え、
$X\mapsto\varepsilon^{-1}$ と $X\mapsto-\varepsilon^{-1}$ は二つの無限遠orderingを
与える。$r$ にsimple rootを持つpolynomialの先頭Laurent項はそれぞれ
$P'(r)\varepsilon$ と $-P'(r)\varepsilon$ である。無限遠では通常のleading coefficientと
$\varepsilon^{-1}$ の対応する冪が先頭項になる。以下の「片側から近づく」「無限遠の符号」
は全てこの構成を意味する。$R((\varepsilon))$ のreal closureへ移しても符号は変わらない。

\begin{lemma}[Legendre congruence]\label{lem:legendre-congruences-ja}
$A,B,C$ を非零、squarefree、pairwise coprime とする。三係数を同じ strict sign にする
ordering がなければ、任意の monic irreducible $\pi$ について
\[
\begin{array}{ll}
 -BC\text{ は modulo }\pi\text{ で平方} &(\pi\mid A),\\
 -AC\text{ は modulo }\pi\text{ で平方} &(\pi\mid B),\\
 -AB\text{ は modulo }\pi\text{ で平方} &(\pi\mid C)
\end{array}                                             \tag{6}
\]
となる。従って同じ三条件はそれぞれ modulo $A,B,C$ でも成り立つ。
\end{lemma}
\begin{proof}
$\pi\mid A$ の場合を示せば、他は cyclic に同じである。pairwise coprime 性から
$BC$ は modulo $\pi$ で非零である。

$\pi=X-r$ が一次とする。$-BC$ が modulo $\pi$ で平方でないなら、実閉体 $R$ では
非零平方が正元と一致するので $-B(r)C(r)<0$、従って $B(r),C(r)$ は同符号である。
$r$ の十分近くでは両者の符号は変わらない。一方、squarefree 性から $A$ は $r$ で
単純根を持ち、左右で符号が反転する。$A$ の符号が $B,C$ と一致する側から $r$ へ
近づく Laurent ordering を取ると、三係数が同符号になり矛盾する。

$\pi$ が非線形なら、実閉体上の既約 polynomial なので $\deg\pi=2$ である。
$R[X]/(\pi)$ は代数閉体 $R(\sqrt{-1})$ であり、その全ての非零元は平方である。
従ってこの場合も (6) が成り立つ。

最後に squarefree polynomial は相異なる既約因子の積である。各因子 modulo の
平方根を中国剰余定理で貼り合わせると、polynomial 全体 modulo の平方根を得る。
\end{proof}

\subsection{Polynomial pivot descent}

次の補題が Lean の ternary proof の停止性を与える。``real split'' は、任意の
実閉順序体への homomorphism で移した symbol が split であることを意味する。

\begin{lemma}[Cyclic pivot step]\label{lem:pivot-ja}
$A,B,C$ を非零squarefree pairwise-coprime polynomials とし、(6) の whole-modulus
conditionsを満たし、$\mathcal Q(A,B,C)$ が real split とする。少なくとも一係数が
非定数なら、$\mathcal Q(A,B,C)$ がすでに $F$ 上 split であるか、または非零squarefree
pairwise-coprime $A',B',C'$ が存在して
\[
\begin{split}
 \mathcal Q(A,B,C)\text{ split}
   &\Longleftrightarrow\mathcal Q(A',B',C')\text{ split},\\
 \mathcal Q(A,B,C)\text{ real split}
   &\Longleftrightarrow\mathcal Q(A',B',C')\text{ real split},\\
 \deg A'+\deg B'+\deg C'
   &<\deg A+\deg B+\deg C.
\end{split}                                             \tag{7}
\]
となり、新しい triple も (6) を満たす。
\end{lemma}
\begin{proof}
まず cyclic permutation により $A$ を pivot に選ぶ。三つの正な次数が全て等しい場合
を除けば、最大次数の係数を選ぶことにより
\[
  \deg B+\deg C<2\deg A                                \tag{8}
\]
とできる。(6) により $x^2\equiv-BC\pmod A$ を満たす $\deg x<\deg A$ の代表を取り
\[
  x^2+BC=Ad                                             \tag{9}
\]
と定める。$x$ の次数boundと (8) から、$d=0$ または $\deg d<\deg A$ である。

$d=0$ なら $x^2=-BC$ であり、$(0,C,x)$ は $\langle A,B,C\rangle$ の
非零 isotropic vector である。これが terminal branch である。

$d\ne0$ とする。cyclic permutation により $\mathcal Q(A,B,C)$ の split 性は
$(-BC,-AB)$ の split 性と同値である。(5) で $u=-BC$、norm を
$x^2+BC=Ad\ne0$ とすると第二 slot は $-A^2Bd$ になる。
平方因子 $A^2$ を除けば $(-BC,-Bd)=\mathcal Q(B,C,d)$ を得る。従って
\[
  \mathcal Q(A,B,C)\text{ split}
  \quad\Longleftrightarrow\quad
  \mathcal Q(B,C,d)\text{ split}                        \tag{10}
\]
である。これらの変換は、分母と norm の非零性が保たれるため、任意の体埋め込み後にも
成立する。

上で説明した parity normalization を triple $(B,C,d)$ に適用する。
各既約因子について偶数valuationを除き、その因子が二つまたは三つのslotに現れる場合は
全slotをその因子倍して新しく生じた平方を除く。出現数は $2$ から $1$、または $3$ から
$0$ へ減る。従ってtotal degreeを増やさずに squarefree pairwise-coprime triple
$(A',B',C')$ を得る。共通scalar倍と各平方置換はsplit性とreal split性を保つので、
(10) から (7) の最初の二つの同値を得る。

新しい triple を同じ strict sign にする ordering は存在しない。もし存在すれば、その
orderingをreal closureへ延長したとき、cleared symbolの二つのslotがともに負になり、
(3) によりnonsplitとなってreal split性に反する。
補題~\ref{lem:legendre-congruences-ja}と中国剰余定理により (6) のwhole-modulus conditionsを
再び得る。Leanではさらに (9) に沿う直接のtransportも記録しており、各coprime pairの
Bézout identityを使って旧square witnessから新しいwitnessを作ってから同じnormalizationを
行う。

次数については、平方・gcd除去前でも新しいtotal degreeは
\[
  \deg B+\deg C+\deg d
  <\deg A+\deg B+\deg C
\]
であり、除去操作は次数を増やさない。従って (7) が成り立つ。

残る場合は $\deg A=\deg B=\deg C=m>0$ である。$+\infty,-\infty$ の ordering は
Laurent series で実現されるので、real split 性から三つのleading coefficientが
同符号になることはない。cyclic permutation後
$\operatorname{lc}(B)\operatorname{lc}(C)<0$ としてよい。実閉性により
\[
  r^2=-\operatorname{lc}(B)\operatorname{lc}(C)
\]
となる $r\in R$ がある。modulo $A$ の reduced square root $x_0$ を
$x=x_0+tA$、$t=r/\operatorname{lc}(A)$ で置き換える。合同は変わらないが、
$x^2+BC$ の次数 $2m$ の係数は相殺する。従って (9) で再び
$d=0$ または $\deg d<m$ となり、上の議論へ帰着する。
\end{proof}

\begin{proposition}[Polynomial quaternion--Legendre criterion]\label{prop:legendre-ja}
$A,B,C\in R[X]$ を非零、squarefree、pairwise coprime とする。
$\langle A,B,C\rangle$ に ordering obstruction がなければ、全て零ではない
$x,y,z\in R[X]$ が存在して
\[
  Ax^2+By^2+Cz^2=0
\]
となる。
\end{proposition}

\begin{proof}
補題~\ref{lem:legendre-congruences-ja}により (6) が成り立つ。任意の実閉imageでは
三係数が同符号でないため、$-AB,-AC$ の少なくとも一方は正である。実閉体の正元は
平方なので、(3) から $\mathcal Q(A,B,C)$ は real split である。

$N=\deg A+\deg B+\deg C$ に関する強帰納法で $F$ 上 split であることを示す。
$N=0$ なら三係数は非零定数で同符号でないから、$-AB,-AC$ の少なくとも一方が正、
従って平方であり、(3) から symbol は split する。$N>0$ では
補題~\ref{lem:pivot-ja}を適用する。terminal branchなら終了し、successor branchなら
total degreeが小さく、(6) と real split 性が保たれる。帰納法の仮定が successor を
splitし、(7) の同値が元の symbol を splitする。

(3) により $\mathcal Q(A,B,C)=(-AB,-AC)$ の split から
\[
  x^2+AB\,y^2+AC\,z^2=0
\]
を満たす非零 $(x,y,z)$ を得る。$A$ で割ると
\[
  A(x/A)^2+B y^2+C z^2=0
\]
である。従って $(x/A,y,z)$ は非零有理関数isotropic vectorであり、共通分母を払えば
全て零でない polynomial solution を得る。
\end{proof}

\begin{proof}[定理~\ref{thm:ternary-ja}の証明]
一方向はscalar extensionである。逆に、全実閉拡大上のisotropyと
補題~\ref{lem:rcdiag-ja}、命題~\ref{prop:orc-ja}から、三係数を同符号にするorderingは
存在しない。上で説明した squarefree pairwise-coprime normalizationを行い、
三項方程式をquaternion symbol
\[
  \left(-\frac ba,-\frac ca\right)
\]
のsplit equationへ移す。命題~\ref{prop:legendre-ja}のpolynomial solutionを得た後、
共通scalar倍と各平方置換を逆にたどれば、元の $\langle a,b,c\rangle$ に対する
$F$ 上の非零isotropic vectorを得る。
\end{proof}

\section{符号分離compression係数}

係数族をbinary head $h=(h_0,h_1)$ と非空tail $t=(t_j)_{j\in J}$ に分ける。

\begin{definition}[Required sign]\label{def:reqsign-ja}
ordered fieldの各点で、補助係数 $s$ のrequired signを
\[
  (h_0>0\land h_1>0)\ \lor\ (\forall j,\ t_j<0)
\]
のとき負、それ以外のとき正と定める。
\end{definition}

\begin{lemma}[Compression sign lemma]\label{lem:compress-sign-ja}
結合族 $h\sqcup t$ がcommon strict signを持たず、$s\ne0$ がrequired signを持つとする。
このとき
\[
  \langle h_0,h_1,s\rangle,
  \qquad
  \langle (t_j)_{j\in J},-s\rangle
\]
のいずれもcommon strict signを持たない。
\end{lemma}
\begin{proof}
headが全正なら $s<0$。headが全負なら、全体は全負でないのでtailは全負でなく、
$s>0$。head が零係数または異符号の係数を含むなら、すでに common strict sign はない。
従って第一形式は common strict sign を持たない。tailが全負なら $s<0$ なので $-s>0$。
tailが全正ならheadは全正であり得ず、$s>0$ なので $-s<0$。tail自身が零係数または異符号を
含む場合は自明である。
\end{proof}

重要なのは、全orderingで同時にこの符号を実現する一つの有理関数を作ることである。

\begin{lemma}[Squarefree polynomial representative]\label{lem:squarefree-normal-ja}
任意の $f\in F^\times$ は
\[
  f=Aq^2
\]
（$A\in R[X]$ は非零squarefree、$q\in F^\times$）と書ける。
\end{lemma}
\begin{proof}
$f=p/r$（$p,r\in R[X]$ は非零）と書けば $f=(pr)/r^2$ である。
unique factorizationにより $pr=Ah^2$ と書く。ここで $A$ は非零の定数単元と、
奇数multiplicityで現れる monic 既約因子の積なのでsquarefreeである。
従って $f=A(h/r)^2$ となる。
\end{proof}

\begin{proposition}[Signed-root skeleton]\label{prop:skeleton-ja}
regular family $h\sqcup t$ にcommon strict-sign orderingがなければ、任意の実閉ordered
extensionでrequired signを持つ $s\in F^\times$ が存在する。
\end{proposition}
\begin{proof}
補題~\ref{lem:squarefree-normal-ja}を用いて各係数を
\[
  a_i=A_iq_i^2,
  \qquad A_i\in R[X]\setminus\{0\}\text{ squarefree}
\]
と書く。平方倍はstrict signとisotropyを変えない。全 $A_i$ の実根全体を有限集合
$\Sigma\subset R$ とする。$R\setminus\Sigma$ の各成分では全係数の符号、従って
required signは一定である。

$r\in\Sigma$ の左右でrequired signが変わる場合に限り $r$ を $E$ に入れる。
一つの外側区間で符号を合わせる $\varepsilon\in\{1,-1\}$ を取り
\[
  S(X)=\varepsilon\prod_{r\in E}(X-r)
\]
と置く。$E$ の根を横切るときだけ $S$ の符号も反転するため、breakpointを順に
横切る帰納により、$S$ は $\Sigma$ 外の全実点でrequired signを持つ。

次に $\varphi:F\to L$ を実閉順序体への準同型、$\xi=\varphi(X)$ とする。
$X-r\ne0$ とinjectivityから $\xi\ne\varphi(r)$。$\Sigma$ に対して $\xi$ と同じ
finite cutを持つ $y\in R$ を選ぶ。隣接する二点の間なら中点、外側の区間なら端点から
1だけ外の点を取ればよい。$\Sigma$ が空なら $y=0$ とする。
実閉順序体間のfield homomorphismは正元を非零平方へ送るので
orderを保つ。実閉体上のpolynomialは実一次因子、根を持たず一定符号の既約二次因子、
leading coefficientの積に分解できる。一次因子の符号はfinite cutだけで決まるから、
$y$ と $\xi$ で全 $A_i$ と $S$ の符号は一致する。ゆえに $\varphi(S)$ がrequired
signを持つ。$s=S\in F$ とすればよい。
\end{proof}

\section{次元帰納}

\begin{lemma}[Binary--tail reconstruction]\label{lem:reconstruct-ja}
$J$ を有限、$s\ne0$ とする。
$\langle a,b,s\rangle$ と
$\langle(\psi_j)_{j\in J},-s\rangle$ がともにisotropicなら、
$\langle a,b,(\psi_j)_{j\in J}\rangle$ もisotropicである。
\end{lemma}
\begin{proof}
第二形式のisotropic vectorを $(y,t)$ とすると
\[
  \sum_j\psi_jy_j^2=s t^2.
\]
$t=0$ ならtail自身がisotropicである。$t\ne0$ なら第一形式のisotropic vector
$(u,v,w)$ を取り、$s=\sum_j\psi_j(y_j/t)^2$ を代入して
\[
  au^2+bv^2+\sum_j\psi_j\left(\frac{wy_j}{t}\right)^2=0
\]
を得る。このvectorが零なら $u=v=0$ となり、$s\ne0$ と第一のisotropyから
$w=0$ となるので矛盾する。
\end{proof}

\begin{proposition}[帰納step]\label{prop:step-ja}
$J\ne\varnothing$ とする。$J\sqcup\{*\}$ 添字の対角形式について
ordering obstructionがないことからisotropyが従うなら、
$\{0,1\}\sqcup J$ 添字についても同じことが従う。
\end{proposition}
\begin{proof}
regular family $h\sqcup t$ にcommon strict-sign orderingがないとする。
命題~\ref{prop:skeleton-ja}の $s$ を取る。任意の実閉ordered extensionで、そのorderを
$F$ へ制限し、元のobstruction仮定と補題~\ref{lem:compress-sign-ja}を使える。
補題~\ref{lem:rcdiag-ja}により二つのcompressed formはともにisotropicである。

三項local--global定理から $\langle h_0,h_1,s\rangle$ は $F$ 上isotropicである。
augmented tail $\langle t,-s\rangle$ を definite にする ordering があれば、
命題~\ref{prop:orc-ja}でその実閉拡大へ移す。この拡大では $s$ が required sign を持ち、
補題~\ref{lem:compress-sign-ja}の第二の結論に反する。従って ordering obstruction はなく、帰納仮定で
$F$ 上isotropicである。
補題~\ref{lem:reconstruct-ja}を適用する。
\end{proof}

\begin{proof}[定理~\ref{thm:diagonal-ja}の証明]
補題~\ref{lem:definite-ja}が (i)$\Rightarrow$(ii) を与える。
(ii)$\Rightarrow$(i) の三次元の場合は、任意の実閉imageで補題~\ref{lem:rcdiag-ja}を
使い、定理~\ref{thm:ternary-ja}を適用すればよい。命題~\ref{prop:step-ja}により
全 $n\ge3$ へ次元帰納する。\texttt{Fin n} 上の結果を有限型とのbijectionでreindexする。
零係数の場合は対応するcoordinate vectorを使う。

(i)$\Rightarrow$(iii) はscalar extensionである。逆にdefinite orderingが存在すると仮定し、
そのordered real closureへ移すと補題~\ref{lem:rcdiag-ja}によりanisotropicとなり、(iii)に
反する。従って(ii)が成り立ち、既に示した含意から(i)を得る。
\end{proof}

\section{一般二次形式}

\begin{lemma}[Orthogonal diagonalization]\label{lem:diagonalization-ja}
標数が $2$ でない体上の有限次元二次形式はorthogonal basisを持つ。
\end{lemma}
\begin{proof}
$B_q(x,y)=q(x+y)-q(x)-q(y)$ をassociated symmetric bilinear formとすると
$B_q(v,v)=2q(v)$ である。次元帰納する。$q=0$ なら任意のbasisでよい。
$q\ne0$ なら $q(v)\ne0$ となる $v$ を選び、各 $x$ に対し
\[
  x^\perp=x-\frac{B_q(x,v)}{2q(v)}v
\]
と置く。$B_q(x^\perp,v)=0$ であり、各 $x$ は $v$ のscalar倍と $v^\perp$ の元の
和として一意に書ける。従って $V=Fv\oplus v^\perp$ である。$v^\perp$ 上の制限に
帰納法の仮定を使い、得られたorthogonal basisへ $v$ を加えればよい。
\end{proof}

\begin{proof}[定理~\ref{thm:qf-ja}の証明]
標数は零なので、補題~\ref{lem:diagonalization-ja}によりorthogonal basis
$e_1,\ldots,e_n$ を取れる。このbasisで $q$ は
\[
  \sum_{i=1}^n q(e_i)X_i^2
\]
とisometricである。退化する場合には一部の $q(e_i)$ が零になり得るが、
定理~\ref{thm:diagonal-ja}は零係数も扱う。isometryはscalar extension後もisometryであり、
invertible linear changeは非零isotropic vectorの存在を保つ。係数族 $(q(e_i))$ に
定理~\ref{thm:diagonal-ja}を適用し、両方向をこれらのisometryで輸送する。
\end{proof}

\section{Lean 宣言との対応}

\begingroup
\scriptsize
\begin{center}
\begin{tabular}{p{0.30\linewidth}p{0.63\linewidth}}
数学上の結果 & Lean 宣言 \\ \hline
Definiteness obstruction & \lean{RatFuncWittLocalGlobal.RatFunc.not_finiteSameStrictSignOrdering_of_isotropic} \\
実閉体上の対角判定 & \lean{RatFuncWittLocalGlobal.Diagonal.not_isotropic_iff_all_pos_or_all_neg} \\
Preordering extension & \lean{RingPreordering.exists_le_isOrdering} \\
正な射影による判定 & \lean{RingPreordering.exists_orderedAlgebraExtension_of_projection} \\
奇数次power-basis lemma & \lean{RatFuncWittLocalGlobal.OrderedRealClosure.minus_one_notMem_baseSquareSpan_of_odd_powerBasis} \\
Ordered real closure & \lean{RatFuncWittLocalGlobal.OrderedRealClosure.orderedFieldRealClosedExtension_same_universe} \\
二次拡大の最大性 & \lean{RatFuncWittLocalGlobal.OrderedRealClosure.nonnegative_isSquare_of_isMax} \\
奇数次拡大の最大性 & \lean{RatFuncWittLocalGlobal.OrderedRealClosure.exists_root_of_odd_natDegree_of_isMax} \\
三項定理 & \lean{RatFuncWittLocalGlobal.RatFunc.ternary_isotropic_iff_forall_realClosedExtension_same_universe} \\
有限Legendre congruence & \lean{RatFuncWittLocalGlobal.RatFunc.PolynomialLegendreLocalConditions} \\
Whole-modulus constructor & \lean{RatFuncWittLocalGlobal.polynomialTernaryFiniteTameResiduesTrivial_globalSquareMod} \\
非線形剰余の平方性 & \lean{RatFuncWittLocalGlobal.RatFunc.nonlinearIrreducibleDivisorSquareMod_of_realClosed} \\
Cleared-binary normalization & \lean{RatFuncWittLocalGlobal.RatFunc.normalizedTernaryLocalGlobal_iff_squarefree_nonterminal_right} \\
Pivot successor & \lean{RatFuncWittLocalGlobal.polynomialTernaryClearedQuaternionSplit_descent_wholeMod_normalized_right_pivot_complete} \\
Total-degree induction & \lean{RatFuncWittLocalGlobal.polynomialTernaryClearedQuaternionSplit_of_wholeMod_realSplit} \\
Legendre endpoint bridge & \lean{RatFuncWittLocalGlobal.polynomialLegendreConstruction_of_normalizedOrderingRealization} \\
Squarefree normal form & \lean{RatFuncWittLocalGlobal.finitePolynomialState_exists} \\
Compression sign lemma & \lean{RatFuncWittLocalGlobal.binaryTailCompressionRequiredSign_spec} \\
Skeleton transport & \lean{RatFuncWittLocalGlobal.FinitePolynomialState.binaryTailSkeleton_allOrdered} \\
Finite-cut sign transport & \lean{RatFuncWittLocalGlobal.polynomial_eval₂_pos_iff_eval_of_same_finite_cut} \\
Reconstruction & \lean{RatFuncWittLocalGlobal.Diagonal.isotropic_sumType_of_ternary_of_tailAdjoinNeg} \\
帰納step & \lean{RatFuncWittLocalGlobal.diagonalOrderingObstructionSufficiency_sum_fin_two} \\
Total diagonal theorem & \lean{RatFuncWittLocalGlobal.diagonal_isotropic_iff_forall_realClosedExtension_of_card_ge_three_same_universe_total} \\
Isometry invariance & \lean{RatFuncWittLocalGlobal.quadraticMap_isometryEquiv_not_anisotropic_iff} \\
Base-change diagonal isometry & \lean{RatFuncWittLocalGlobal.quadraticFormBaseChangeIsometryWeightedSumSquares} \\
一般二次形式 & \lean{RatFuncWittLocalGlobal.quadraticForm_not_anisotropic_iff_forall_realClosedExtension_same_universe}
\end{tabular}
\end{center}
\endgroup

再現可能な公理監査 \texttt{scripts/AxiomAudit.lean} は、主定理が依存する公理が
\[
  \texttt{propext},\qquad \texttt{Classical.choice},\qquad \texttt{Quot.sound}
\]
のみであることを検査し、異なる場合には失敗する。project固有のaxiomはない。
証明項に現れる bundle と transport の宣言は、本文で証明した
合同条件、正規化、finite cut、座標変換を Lean の型に実装するものであり、本文にない
追加の数学的定理を仮定するものではない。

\begin{thebibliography}{9}
\bibitem{Lam}
T. Y. Lam,
\emph{Introduction to Quadratic Forms over Fields},
Graduate Studies in Mathematics 67, American Mathematical Society, 2005.

\bibitem{Pfister}
A. Pfister,
\emph{Quadratic Forms with Applications to Algebraic Geometry and Topology},
London Mathematical Society Lecture Note Series 217,
Cambridge University Press, 1995.
\end{thebibliography}

\end{document}
